\documentclass[11pt]{article}

\usepackage[margin=1in]{geometry}
\usepackage{amsmath,amssymb,amsthm}
\usepackage{hyperref}
\usepackage{graphicx}

\newtheorem{theorem}{Theorem}
\newtheorem{remark}{Remark}

\newcommand{\BPBp}{\operatorname{BPBp}}
\newcommand{\RSE}{\operatorname{RSE}}
\newcommand{\SE}{\operatorname{SE}}
\newcommand{\str}{\operatorname{str}}
\newcommand{\dist}{\operatorname{dist}}
\newcommand{\norm}[1]{\left\lVert #1 \right\rVert}

\title{A Hilbert-Range Subcase of the Finite-Dimensional Domain BPBp-RSE Question}
\author{FA/Banach discovery run \texttt{fa\_banach\_001}}
\date{2026-06-14}

\begin{document}
\maketitle

\section*{Source Question}

The source paper is Helena del Rio, \emph{Bollobas-type theorems for range
strongly exposing operators}, arXiv:2512.10442.  It ends with the following
question.

\begin{quote}
Suppose that \(X\) has finite dimension and \(Y\) is uniformly convex.
Does the pair \((X,Y)\) have the BPBp(-RSE)?
\end{quote}

\begin{center}
\includegraphics[width=0.86\linewidth]{figures/open_problem_crop.png}
\end{center}

This packet gives a positive answer for the Hilbert-range subcase.  It is
not a solution of the full uniformly convex range question.

\section*{Definitions Used}

A pair \((X,Y)\) has the Bishop-Phelps-Bollobas property for operators
(\(\BPBp\)) if, for every \(\varepsilon>0\), there is
\(\eta(\varepsilon)>0\) such that whenever \(T\in \mathcal L(X,Y)\),
\(\norm{T}=1\), and \(x\in S_X\) satisfy
\[
  \norm{Tx}>1-\eta(\varepsilon),
\]
there are \(S\in\mathcal L(X,Y)\) and \(u\in S_X\) such that
\[
  \norm{S}=\norm{Su}=1,\qquad
  \norm{S-T}<\varepsilon,\qquad
  \norm{u-x}<\varepsilon.
\]

The source paper defines \(\BPBp\)-\(\RSE\) by requiring, in addition, that
the approximating operator \(S\) be range strongly exposing.  We use
Corollary 2.8 of the source paper: if a subspace of \(\mathcal L(X,Y)\)
containing the rank-one operators has the classical \(\BPBp\), and
\(\str(B_Y)=S_Y\), then it has the \(\BPBp\)-\(\RSE\).  Hilbert spaces are
uniformly convex, hence every point of their unit sphere is strongly exposed.

\section*{Proof Intuition}

The only point where the infinite-dimensional Hilbert range could cause a
loss of compactness is the location of the finite-dimensional range
\(T(X)\).  But Hilbert geometry removes this issue: every \(m\)-dimensional
Hilbert subspace is isometric to the same model space \(\ell_2^m\).  If
\(\dim X=n\), then an operator \(X\to H\) only sees one of the finitely many
Hilbert dimensions \(1,\ldots,n\).  Taking the minimum of the known
finite-dimensional BPBp moduli over those finitely many models gives a
single modulus for the whole Hilbert space.

\section*{Theorem}

\begin{theorem}
Let \(X\) be a finite-dimensional real or complex Banach space and let \(H\)
be a Hilbert space over the same scalar field.  Then \((X,H)\) has the
\(\BPBp\)-\(\RSE\).
\end{theorem}

\begin{proof}
Let \(n=\dim X\).  We first prove the classical \(\BPBp\) for \((X,H)\).

Fix \(0<\varepsilon<1\).  For each \(1\leq m\leq n\), let \(H_m\) denote an
\(m\)-dimensional Hilbert space.  Since both \(X\) and \(H_m\) are
finite-dimensional, the finite-dimensional Bishop-Phelps-Bollobas theorem
for operators of Acosta--Aron--Garcia--Maestre gives
\[
  \eta_m=\eta(X,H_m)(\varepsilon)>0 .
\]
Put
\[
  \eta=\min_{1\leq m\leq n}\eta_m>0 .
\]

Let \(T\in \mathcal L(X,H)\) satisfy \(\norm{T}=1\), and let \(x\in S_X\)
satisfy \(\norm{Tx}>1-\eta\).  Set \(E=T(X)\).  Then \(E\) is a Hilbert
subspace of \(H\) of some dimension \(1\leq m\leq n\).  Choose a surjective
linear isometry \(U:E\to H_m\).  The operator \(UT:X\to H_m\) has norm one
and
\[
  \norm{UTx}=\norm{Tx}>1-\eta\geq 1-\eta_m .
\]
By the \(\BPBp\) for \((X,H_m)\), there are \(R\in\mathcal L(X,H_m)\) and
\(u\in S_X\) such that
\[
  \norm{R}=\norm{Ru}=1,\qquad
  \norm{R-UT}<\varepsilon,\qquad
  \norm{u-x}<\varepsilon .
\]
Define \(S=U^{-1}R:X\to E\subset H\).  Since \(U^{-1}\) is an isometry,
\[
  \norm{S}=\norm{Su}=1,\qquad
  \norm{S-T}=\norm{R-UT}<\varepsilon,\qquad
  \norm{u-x}<\varepsilon .
\]
Thus \((X,H)\) has the classical \(\BPBp\).

Finally, Hilbert spaces are uniformly convex.  Hence \(\str(B_H)=S_H\).
Applying Corollary 2.8 of the source paper to the full space
\(\mathcal L(X,H)\), which contains all rank-one operators, upgrades the
classical \(\BPBp\) to the \(\BPBp\)-\(\RSE\).
\end{proof}

\section*{Limitations and Novelty Check}

This is only a Hilbert-range subcase of Question 5.11.  The proof uses the
fact that all finite-dimensional Hilbert subspaces of a fixed dimension are
isometric.  For a general uniformly convex range \(Y\), the subspaces
\(T(X)\) of dimension at most \(\dim X\) can have many non-isometric norms.

A tempting compactness argument over finite-dimensional range norms is not
valid by itself: the BPBp modulus can degenerate along finite-dimensional
Banach-Mazur convergent ranges.  The ACKLM example of two-dimensional
polyhedral spaces \(Y_k\) with
\(\inf_k \eta(\ell_1^2,Y_k)(\varepsilon)=0\) shows that Banach-Mazur
compactness alone cannot give the desired uniform lower bound.  A full proof
for uniformly convex ranges would need a genuinely uniform-convexity-based
stability argument, not merely finite-dimensional compactness.

Bounded literature checks on 2026-06-14 included exact and close searches
for finite-dimensional domain BPBp with uniformly convex range, the source
paper's final question, the ACKLM universal BPB-domain paper
arXiv:1305.6420, and the finite-dimensional approach papers for
\(\ell_\infty^3\) and \(\ell_\infty^4\).  These checks found known positive
results for finite-dimensional ranges, \(\ell_1^m\) domains, \(\ell_\infty^n\)
domains with uniformly convex ranges, and \(C(K)\) domains with uniformly
convex ranges, but no full answer to Question 5.11.  The Hilbert subcase may
be regarded as an elementary consequence of the finite-dimensional BPBp
theorem rather than a deep new theorem.

\section*{Human Review Recommendation}

Check that the finite minimum of the finite-dimensional moduli is used only
over Hilbert isometry classes \(H_m\), \(1\le m\le \dim X\), and that the
RSE upgrade is exactly Corollary 2.8 of the source paper.  The packet should
not be read as evidence that the arbitrary uniformly convex range case is
settled.

\begin{thebibliography}{9}

\bibitem{AAGM2008}
M. D. Acosta, R. Aron, D. Garcia and M. Maestre,
\emph{The Bishop-Phelps-Bollobas theorem for operators},
J. Funct. Anal. 254 (2008), 2780--2799.

\bibitem{ACKLM2015}
R. Aron, Y. S. Choi, S. K. Kim, H. J. Lee and M. Martin,
\emph{The Bishop-Phelps-Bollobas version of Lindenstrauss properties A and B},
Trans. Amer. Math. Soc. 367 (2015), 6085--6101; arXiv:1305.6420.

\bibitem{delRio2025}
H. del Rio,
\emph{Bollobas-type theorems for range strongly exposing operators},
arXiv:2512.10442.

\end{thebibliography}

\end{document}
